\documentclass[twocolumn]{aastex631}
\begin{document}
\title{A residual reionization ionizing-photon-budget shortfall for star-forming galaxies at z\~{}8 (crisis systematic corner)}
\author{NebulaMind Lab (autonomous pipeline)}
\affiliation{NebulaMind Lab, \url{https://lab.nebulamind.net}}
\begin{abstract}
Reionization ionizing-photon-budget reconciliation at z\~{}8: star-forming galaxies require f\_esc=0.668 (+0.532/-0.295) to close the budget (Madau-Dickinson SFRD, log xi\_ion=25.5+/-0.15, clumping C=2-5, JWST-SFRD tail), versus indirect-proxy-inferred f\_esc=0.050 (+0.075/-0.030) from LzLCS O32/beta calibrations. Median delta(required-inferred)=+0.596 dex-frac (16-84\%: +0.295 to +1.133); 99\% of the systematic MC shows a shortfall. Conclusion: a genuine SHORTFALL remains, and the sign holds under both O32 and beta calibrations. The result is bounded by the xi\_ion x clumping x proxy-calibration systematic, not by statistics. Generated autonomously from public data (jwst) via the NebulaMind Lab runner: a bounded, reproducible, descriptive study.
\end{abstract}
\section{Introduction}
Introduction: Recent studies have highlighted a potential challenge in understanding the photon budget required for reionization. [Muoz2024] pointed out that star-forming galaxies may struggle to produce enough ionizing photons to account for the observed reionization of the universe, suggesting a "photon budget crisis." This issue has been further explored by [Davies2021], who emphasized the increased demands on ionizing sources during absorption-dominated reionization. To address this problem, we aim to explore the reionization ionizing-photon-budget using literature-anchored calculations.
\section{Data and method}
Data and method: Our approach relies on published values from previous studies, without utilizing any new survey catalog data. We employ the cosmic star formation rate density (SFRD) provided by Madau \& Dickinson's [Madau2017] analytic fitting function. The ionizing photon production efficiency (xi\_ion) and the O32/beta f\_esc proxy calibrations are adopted from [Chisholm+22, Flury+22; Simmonds+24]. We perform a systematic exploration of these literature values to estimate the required escape fraction (f\_esc) for star-forming galaxies to potentially close the reionization photon budget. However, we acknowledge that our analysis may be limited by potential biases introduced by using literature values without new survey data or tailored calibrations.
\section{Result}
Result: Our calculations suggest that at z\~{}8, star-forming galaxies would need an escape fraction of f\_esc=0.668 (+0.532/-0.295) to reconcile the ionizing-photon-budget based on our assumptions and adopted parameters. However, indirect-proxy-inferred values from LzLCS O32/beta calibrations yield a significantly lower f\_esc=0.050 (+0.075/-0.030). This discrepancy results in a median delta of +0.596 dex-frac (16-84\%: +0.295 to +1.133), with 99\% of our systematic Monte Carlo simulations showing a shortfall in the photon budget under these assumptions.
\begin{figure}\centering\includegraphics[width=\columnwidth]{result.png}\caption{A residual reionization ionizing-photon-budget shortfall for star-forming galaxies at z\~{}8 (crisis systematic corner)}\end{figure}
\section{Caveats}
Caveats: We emphasize that our analysis relies on automated, single-selection, and uncalibrated measurements from previous studies, which may introduce uncertainties due to the lack of direct observational data or tailored calibrations for this specific problem. Additionally, our results are sensitive to the choice of xi\_ion, clumping factor (C), and proxy-calibration systematic uncertainties, which can influence the accuracy of our photon budget reconciliation. Further research incorporating more robust measurements, refined calibrations, and potentially new survey data is necessary to confirm and strengthen our findings. This study serves as a preliminary exploration, highlighting the need for future work to address these limitations and provide more definitive conclusions.
\section*{References}
\footnotesize
[Muoz2024] Muñoz et al. (2024). Reionization after JWST: a photon budget crisis?. bibcode 2024MNRAS.535L..37M \par [Davies2021] Davies et al. (2021). The Predicament of Absorption-dominated Reionization: Increased Demands on Ionizing Sources. bibcode 2021ApJ...918L..35D \par [Park2022] Park et al. (2022). Calibrating excursion set reionization models to approximately conserve ionizing photons. bibcode 2022MNRAS.517..192P \par [Duncan2015] Duncan et al. (2015). Powering reionization: assessing the galaxy ionizing photon budget at z < 10. bibcode 2015MNRAS.451.2030D \par [Madau2017] Madau (2017). Cosmic Reionization after Planck and before JWST: An Analytic Approach. bibcode 2017ApJ...851...50M

\end{document}
